% Summary of polinsky2001: Long-distance agreement and topic in {Tsez}

\begin{abstract}
\begin{enumerate}
  \item In Tsez, verbs agree with absolutive arguments in noun class.
  \item When the absolutive argument is a clause, and that clause has an absolutive topic (covertly moved at LF), verbs agree with the absolutive topic. If there is no absolutive topic, they agree with the clause itself.
\end{enumerate}
\end{abstract}

\section*{Data}

\begin{enumerate}
  \item mother [boy bread.III.ABS ate].IV \textbf{IV}-know (PLA)
  \item mother [boy bread.III.ABS ate].IV \textbf{III}-know (LDA)
\end{enumerate}
`The mother knows the boy ate bread.'
